\centerline{\Huge \bf  Олимпиада I}
\centerline{\Huge \bf 8 класс}

\begin{flushright}
    \scriptsize{Заключительный отбор в Сириус 2021г.\\ Кто загуглит - тому бан!}
\end{flushright}

\problem{1}
Неотрицательные числа записаны в клетках квадрата $7$ на $7$. Сумма чисел в любых двух соседних столбцах не менее $20$, а в любых двух соседних строках не более $15$. Чему может быть равна сумма всех чисел, записанных в клетках данного квадрата?

\problem{2}
Дан параллелограмм $ABCD$. На стороне $AD$ выбрана точка $P$ так, что $BP$ является биссектрисой угла $B$. Найдите длины сторон параллелограмма $ABCD$, если $BP = CP = 6 \ \text{и} \ PD = 5.$

\problem{3}
Для положительных чисел $x, y$ таких, что $xy = 4$ докажите неравенство
\[ \dfrac{1}{x+3} + \dfrac{1}{y+3} \leq \dfrac{2}{5}\]

\problem{4}
Найдите все положительные целые решения уравнения
\[(a+1)^4 \cdot (b+1)^4 \cdot (c+1)^4 =(40a+1)\cdot(40b+1)\cdot(40c+1).\]

\problem{5}
Квадрат со стороной 3 и прямоугольник с целыми сторонами, одна из которых на 2 длинее другой, расположены так, что их стороны параллельны, а центры совпадают. Стороны квадрата продлили так, как показано на рисунке. Может ли площадь заштрихованной области быть чётным числом?

\problem{6}
Две криптовалюты были созданы 22 мая 2021 года. Каждая стоит сначала по 1 рублю за монету. Стоимость первой валюты увеличивается на 1 рубль в день. А курс второй валюты каждый день вычисляется как отношение суммы курсов двух валют за прошлый день к произведению этих же двух значений. Найдите курс второй валюты через 3650 дней.